\clearpage
\refstepcounter{chapter}
\phantomsection
\addcontentsline{toc}{chapter}{附录 \thechapter\hspace{0.5em}补充结果}
\chapter*{附录 \thechapter\hspace{1em}补充结果}
\markboth{附录 \thechapter\hspace{1em}补充结果}{}
\label{chap:appendix}

本附录给出补充的指标说明与标准化公式，用于演示附录中表、公式的编号方式（表~\ref{tab:appendix_indicators}、式~\eqref{eq:appendix_norm}）。

\section{指标说明}

表~\ref{tab:appendix_indicators} 列出数字经济发展水平评价的部分基础指标及其单位。

\begin{table}[!htbp]
    \centering
    \caption{数字经济发展水平评价部分基础指标（示例）}
    \label{tab:appendix_indicators}
    \wuhao\songti
    \begin{tabular}{llc}
        \toprule
        维度 & 指标 & 单位 \\
        \midrule
        数字基础设施 & 每百人互联网宽带接入端口数 & 个 \\
        数字产业化   & 信息传输、软件业增加值占比 & \% \\
        产业数字化   & 电子商务销售额 & 亿元 \\
        \bottomrule
    \end{tabular}

    \dufesource{资料来源：本文整理（示例）。}
\end{table}

\section{标准化方法}

各基础指标采用极差法进行标准化，正向指标的标准化公式为

\begin{equation}
    x'_{ij} = \frac{x_{ij}-\min_i x_{ij}}{\max_i x_{ij}-\min_i x_{ij}}.
    \label{eq:appendix_norm}
\end{equation}
